\input amstex
\documentstyle{amsppt}
\magnification 1200
\pageno 43
\define\eps{\epsilon}
\define\okk{{\Cal O}_K}
\define\off{\overline{F}}
\define\ovf{\overline{F}}
\define\unb{\underline{B}}
\define\unt{\underline{T}}
\define\unx{\underline{x}}
\define\qqq{\Bbb Q}
\define\ccc{\Bbb C}
\define\zzz{\Bbb Z}
\define\rrr{\Bbb R}
\define\aaa{\Bbb A}
\define\nnn{\Bbb N}
\define\gga{\Bbb G}
\define\oee{{\Cal O}_E}
\define\VV{\Cal V}
\define\HH{\Cal H}
\hsize 6.5 true in
\vsize 8.7 true in
\vskip .3in
\subheading{\S 6 The matching for the $1$-regular germ}
\vskip .1in
Let $G=G_2$ throughout this section.

The $\kappa$-stable subregular germ for $G_2$ should match with the 
stable subregular germ of $PGL(3)$ or $(SL(2)\times SL(2))/\{\pm 1\}$, 
for homogeniety reasons. The idea to verify this is simply to write down
the Igusa expansions of the two orbital integrals and match the $\beta =2$ 
terms. Furthermore, for the 1-regular comparisons, we can and do
assume that $f$ is supported off the 0-regular orbit.
\par
Let $u=exp(X_\alpha+X_{\alpha+\beta})$. Then the subregular orbit
is given by ${\Cal O}_1=u^{G_2}$. This is not a stable orbit,
but it is Richardson.
See the appendix for more on its $F$-classes.
It will turn out that both the biquadratic and the cubic Cartans
can have non-zero 1-regular $\kappa$-stable germs.
\par
From \S\S 4,5, there are three base divisors with $\beta =2$,
and from \cite{H1, ch. VI, \S 2}, we know that all fundamental divisors
of $G_2$ have $\theta =1$.
They are described as follows: 
\par
$E_\alpha$ is given locally as $\{\lambda =0\}$ 
(with multiplicity one) on $Z_\alpha$
and $Z_\alpha$ is regular with generators of its local coordinate ring
being $\lambda$, $x(\alpha)$, $z(\beta)$, $x(\gamma)$,
$x(\delta)$, $x(\epsilon)$, $x(\zeta)$.
\par
$E_\beta$ is given locally as $\{\lambda =0\}$ (with multiplicity one) 
on $Z_\beta$
and $Z_\beta$ is regular with generators of its local coordinate ring
being $\lambda$, $z(\alpha)$, $x(\beta)$, $x(\gamma)$,
$x(\delta)$, $x(\epsilon)$, $x(\zeta)$.
\par
$E_{\emptyset,\alpha,\gamma}$ is given locally
by $\{z(\alpha) =0\}$ (with multiplicity one) 
on $Z_{\emptyset,\alpha,\gamma}$ and $Z_{\emptyset,\alpha,\gamma}$
is regular with generators of its local coordinate ring
being $z(\alpha)$, $z(\beta)$, $x(\gamma)$, $w(\gamma)$,
$x(\delta)$, $x(\epsilon)$, $x(\zeta)$. This divisor may be eliminated
from consideration by \cite{H1, Ch V, Lemma 1.1}.

As a corollary to \cite{H1, ch 6, Lemma 1.1}, we know that
$E_\alpha$ and $E_\beta$ are $F$-rational. Note that the action of the Galois
group $\Gamma$ on $T_G$ and on $Z$ factors through a finite group
$\Gamma_{T}\subset \Omega$.
To satisfy the hypotheses (5.1-4), one may use
the calculations of \S 2 to ``average'' the coordinates of $Z_\alpha$ or
$Z_\beta$ over the Galois group $\Gamma_T$
to get $F$-rational coordinates on these divisors.
\par
Moreover, if we choose the Cartan $T_G$ such that $\Gamma$ acts
on $T_G$ via a product of two cyclic subgroups of order two of the Weyl group
$\Omega$ then the $\kappa$-stable subregular germ 
of $G_2$ should match with the stable subregular germ of
$(SL(2)\times SL(2))/\{\pm 1\}$ 
(and by the appendix, $\kappa$ is quadratic)
if we choose the Cartan $T_G$ such that $\Gamma$ acts on $T_G$
via a cyclic subgroup of $\Omega$ of order three
then the $\kappa$-stable subregular germ 
of $G_2$ should match with the stable subregular germ of
$PGL(3)$ (and by the appendix, $\kappa$ is a cubic character).
\proclaim{Theorem 6.1} Suppose $T_G$ is a non-split Cartan subgroup
of $G_2$ which splits over a
Galois extension for which the Galois group action on $T_G$
is either a direct product of two cyclic subgroups of $\Omega$ order $2$
or a cyclic subgroup of order $3$.
Assume that $f\in C_c^\infty(G(F))$ satisfies $supp(f)\cap {\Cal O}_1
\not= \emptyset$, where ${\Cal O}_1$ denotes a fixed $1$-regular
$G(F)$-orbit, and the support of $f$ does not meet any unipotent orbit
of $G(F)$ which is not $1$-regular. Then there is an $f^H$, smooth on
$H(F)$, for which, for all regular $\gamma_G\in T_G(F)$
in a sufficiently small neighborhood of $1$, we have
$$
\Delta_{G/H}^{T_G,\kappa}(\gamma_H,\gamma_G)
\Phi^{T_G,\kappa}(\gamma_G,f)=
\Phi^{T_H,st}(\gamma_H,f^H),
$$
where $\kappa$, $T_G$ satisfy
$H_{T_G,\kappa}=H_i$. Here $i=1$, so $H=PGL(3)$, if
the Galois group action on $T_G$
is a cyclic subgroup of order $3$, and
$i=2$, so $H=(SL(2)\times SL(2))/\{\pm 1\}$, if
the Galois group action on $T_G$ is
a direct product of two cyclic subgroups of $\Omega$ order $2$.
\endproclaim
This will be proven below.
\vskip .1in 
{\bf Proof of Theorem 6.1}
\vskip .1in
By the work of Langlands and Shelstad, \cite{LS1}, \cite{LS2}, we need only
prove the corresponding matching of germs conjecture stated in
\S 0 above.

We discuss the Igusa expansion of the $G_2$ side more carefully:
Since only the $\theta=1$ term of formula 0.6(a)
contributes by \cite{H1, ch 6, \S 2}, 
the $\beta=2$ term of the Igusa expansion takes the form, 
$|\lambda|\kappa(inv(\gamma_G^0,\gamma_{G,\lambda}))
D_H(\gamma_{H,\lambda})^{-1}$ 
times
$$
F_1(1,2,f_\kappa)=
PV\int_{U\cap E(F)}A_1(M)h_1|\nu_1|=
PV\int_{E(F)\cap U}\pi^*(f)\kappa_1(t_*)|\omega_E|
\tag 6.2
$$
where $E$ is $E_\alpha$ or $E_\beta$, depending on $T_G$, and
where $A_1(M)$, $h_1$, $\nu_1$, ..., are as in \cite{L},
and $f_\kappa$ includes the $\kappa$-factor. The factor 
$m_\kappa(e)/\kappa_1(\lambda)|_{E(u)}$ is given by
\cite{H1, Ch VI, \S 2 and Proposition 3.2}, provided $E$ is the divisor
associated to a line of the Dynkin curve associated to the long root
of $G_2$ which is fixed by the Galois action. 
By the homogeniety of germs, formula 0.6 and lemma 0.8,
this factor is equal to a multiple of
$$
\sum_{{\Cal O}\subset {\Cal O}_1^{st}}
\mu_{\Cal O} (f)\Gamma^{T_G,\kappa}_{\Cal O}(\gamma_\lambda),
$$
where ${\Cal O}_1^{st}$ denotes the stable 1-regular unipotent orbit.
M. Assem has computed the number of $F$-rational classes in the
stable subregular class when $p\not= 2,3$ (see the appendix below).
\par
We discussed the geometry of the fibering $E_\alpha\rightarrow 
\overline{{\Cal O}}_1$, where $\overline{{\Cal O}}_1$ denotes the
Zariski closure, in \S 5. 
We assume that the function $f$ meets ${\Cal O}_1^{st}$ and no other 
unipotent orbit.
By formula 0.6 and (6.2), we have 
$$
PV\int_{U\cap E(F)}A_1(M)h_1|\nu_1|
=\int_{\pi(U)\cap {\Cal O}_1(F)}f\int_{E_u(F)}
\kappa_1(t_*)|\omega_{E_\alpha(u)}||\omega_{{\Cal O}_1}|,
\tag 6.3
$$
where $E$ denotes either $E_\alpha$ or $E_\beta$.
We may restrict ourselves to 
the case where $T_G$ is either biquadratic or cubic
as in Theorem 6.1. The dependence of (6.3) on 
$C\subset exp(\off\cdot X)$, where $X\in {\frak t}$ is a tangent vector
for $C$, is via the dependence of $\kappa(t_*)$ and $\omega_{E_\alpha(u)}$
on $X$, where ${\frak t}$ denotes the Lie algebra of $T_G$.
\par
Recall \cite{L, Lemma 2.12} that the measure $\omega_Y$ already incorporates 
the Weyl discriminant $\prod_{\sigma >0}(1-\sigma(\gamma)^{-1})$.
Also, from \S 4, recall 
$$
|\omega_{Z_\alpha}|=
{\lambda\over z(\beta)^2}d\lambda dx(\alpha) dz(\beta)dx(\gamma)
dx(\delta)dx(\epsilon)dx(\zeta),
$$
and
$$
|\omega_{Z_\beta}|=
{\lambda\over z(\alpha)^2}d\lambda dz(\alpha) dx(\beta)dx(\gamma)
dx(\delta)dx(\epsilon)dx(\zeta).
$$
Note that these coordinates are not all $F$-rational if $c_\sigma$ is a
simple reflection by (2.16), (2.17).

On ${\Cal O}_1$ the roots ``outside $P_\alpha$'' (or the
roots ``outside $P_\beta$'') may be used to
parameterize the orbit - see the lemma below.

The digression below is motivated by \cite{H1, Remark V.2.2}.
\par
Let ${\Cal O}$ be the Richardson orbit associated to a parabolic $P$
as in \cite{St, Theorem 2,\S 3.9}. Write $P=MN_P$, so that $N_P^G$
contains ${\Cal O}$ as a dense subset, ${\Cal O}\cap N_P$ is dense in
$N_P$ and ${\Cal O}\cap N_P$ forms a single class under $P$. We have
dim${\Cal O}$ = dim$G$ - dim$M$. Let $u\in {\Cal O}\cap N_P$ and recall
$N_P^- P$ is dense in $G$ by the Bruhat decomposition. The $P$-orbit
$\{u^p\ |\ p\in P\}$ is equal to ${\Cal O}\cap N_P$ and is dense in $N_P$.
Therefore, the coordinates of $N_P$ yield coordinates for the $P$-orbit
of $u\in {\Cal O}\cap N_P$. The $G$-orbit of this $u$, that is
${\Cal O}$ itself, contains the $N_P^- P$-orbit of $u$ as an open
subvariety. We may regard the $N_P^- P$-orbit of $u$ as the
$N_P^-$-orbit of ${\Cal O}\cap N_P$. By Corollary 1.6, for each
$\overline{n}\in N_P^-$, $B^{\overline{n}}$ is opposite to
$B^-$ and $B^{\overline{n}}\not= B$ unless $\overline{n}=1$.
Therefore, for our $u$ as above, $u^{\overline{n}}\notin {\Cal O}\cap N_P$.
From this it follows that the $N_P^-$-orbit of ${\Cal O}\cap N_P$
is isomorphic to $N_P^-$.
\proclaim{Lemma 6.4 (Hales)}
If ${\Cal O}$ is the Richardson orbit associated to a
parabolic $P$ then the coordinates on $N_P^- \times N_P$ serve as local
coordinates on ${\Cal O}$.
\endproclaim

Returning to the proof of Theorem 6.1,
associated to the roots outside $P_\alpha$, we have the invariant measure
$$
\omega_{{\Cal O}_1}=
\lim_{\lambda\rightarrow 0}
{dx(\beta)...dx(\zeta)dv(\beta)...dv(\zeta)\over \lambda},
$$
Recall that the invariant measure $\omega_{{\Cal O}_1}$ is
unique up to a constant \cite{R}. On the fiber $E_\beta(u)$, we use
$$
\omega_{E_\beta(u)}={1\over z(\alpha)^2}dv(\alpha)dz(\alpha),
$$
by
\S 4 above, \cite{L, Lemma 2.12}. We have then
$$
\omega_{E_\beta(u)}=
{(1-\alpha (\gamma_\lambda)^{-1})\over \lambda}|_{\lambda=0}
dv(\alpha){dz_\alpha\over z_\alpha^2}
={\alpha(X)}dv(\alpha){dz_\alpha\over z_\alpha^2},
$$
where $\alpha(X):={1-\alpha^{-1}(\gamma_\lambda)\over\lambda}|_{\lambda=0}$ 
denotes the inner product of $\alpha$ with the tangent vector $X$ of $C$.
Similarly, associated to the roots outside $P_\beta$, we have
$$
\omega_{E_\alpha(u)}={1\over z(\beta)^2}dv(\beta)dz(\beta)
={\beta(X)}dv(\beta){dz_\beta\over z_\beta^2}.
$$
Therefore, the $\beta =2$ term is equal to
$$
\split
&|\lambda|\kappa(inv(\gamma_G^0,\gamma_{G,\lambda}))
D_H(\gamma_{H,\lambda})^{-1} 
PV\int_{U\cap (E_\alpha)(F)}A_1(M)h_1|\nu_1|\\
&=
|\lambda|\kappa(inv(\gamma_G^0,\gamma_{G,\lambda}))
D_H(\gamma_{H,\lambda})^{-1} 
|\beta(X)|\int_{{\Cal O}_1}f\int_{E_\alpha(u)\cap (U\cap \pi^{-1}(u))}
\kappa(t_*)|dv'(\beta)
{dz_\beta\over z_\beta^2}||\omega_{{\Cal O}_1}|,
\endsplit
\tag 6.5a
$$
or
$$
\split
&|\lambda|\kappa(inv(\gamma_G^0,\gamma_{G,\lambda}))
D_H(\gamma_{H,\lambda})^{-1} 
PV\int_{U\cap (E_\beta)(F)}A_1(M)h_1|\nu_1|\\
&=|\alpha(X)|\int_{{\Cal O}_1}f\int_{E_\beta(u)\cap (U\cap \pi^{-1}(u))}
\kappa(t_*)|dv'(\alpha)
{dz_\alpha\over z_\alpha^2}||\omega_{{\Cal O}_1}|,
\endsplit
\tag 6.5b
$$
From these expressions and the remarks in \S 0 above, 
one may verify the local translation invariance 
property of Hales' lemma 5.19.
From this, the transfer follows in the $1$-regular case,
for compactly supported functions whose support meets the 1-regular
orbit and no other orbit.
\enddocument
\par
For the sake of comparison, we look at
the corresponding expression for the subregular 
term for $PGL(3)$, $(SL(2)\times SL(2))/\{\pm 1\}$. 
\par
The stable subregular germ of $(SL(2)\times SL(2))/\{\pm 1\}$: 
The stable subregular
unipotent orbit of $(SL(2)\times SL(2))/\{\pm 1\}$ is the union of
${\Cal O}_{tr}\times {\Cal O}_{reg}^{st}$ and
${\Cal O}_{reg}\times {\Cal O}_{tr}$. The stable regular germ on $SL(2)$
may be identified with the regular germ of $PGL(2)$ (they differ
at most by some non-zero constant of normalization -
here one may use another result of Repka or the fact that
$\sum_{{\Cal O}\subset {\Cal O}_H^{st}}C_{\Cal O}(\gamma_\lambda)$
is a constant, in the notation of \cite{SS, Lemma 2.4}, 
where $p\not= 2$ is used) and
can be found, for example, in \cite{SS}; to get the stable 
subregular germ of $(SL(2)\times SL(2))/\{\pm 1\}$ we take the product of
the regular germ of $PGL(2)$ with the stable trivial germ of $SL(2)$. 
On each fixed 
elliptic Cartan of $PGL(2)$ the regular germ is a non-zero constant
times $|D(\gamma_\lambda)|^{-1/2}$, where $D$ is the usual Jacobian.
Denoting the simple positive roots on $(SL(2)\times SL(2))/\{\pm 1\}$
by $\alpha_1$, $\alpha_2$, we find that
$$
|D_{H_0}(\gamma_{\lambda,H})|^{1/2}=
|\lambda||\alpha_1(X_{H_0})|, H_0=SL(2),
\tag 6.6
$$
for $\gamma_{\lambda,H}=exp(\lambda X_H)\in T^H(F)\rightarrow
SL(2,F)\times SL(2,F)$ elliptic regular. We want this function of $X_H$
to be equal, up to a constant factor, to the function of $X$ occuring
in $c_{\kappa, T,C}$. Recall that $\alpha_1$ corresponds to $\alpha$ and
$\alpha_2$ corresponds to $3\alpha +2\beta$.
\par
The stable subregular germ of $PGL(3)$: The stable subregular
unipotent orbit of 
\newline
$PGL(3)$ consists of only one $F$-class. The subregular germ on $GL(3)$
can be found in \cite{Rep}. On  regular elements of an
elliptic Cartan $E^{\times}$ of $GL(3)$, it is a non-zero constant
times $|\lambda (3\alpha +2\beta)(X)||D_H(\gamma_\lambda)|^{-1/2}$
(at least for $p\not= 3$). The matching of the Igusa transfer factor
$p_{{\Cal O},T}$ with the LS transfer factor predicts something slightly 
different: a non-zero constant times
$|\lambda \beta(X)||D_H(\gamma_\lambda)|^{-1/2}$. However, these are 
actually the same quantity, the discrepancy being accounted for by 
the fact that for such elliptic Cartans of $PGL(3)$,
$|1-\beta(\gamma_G)|=|1-(3\alpha +2\beta)(\gamma_G)|$ 
for all regular $\gamma_H\in T_H(F)$.
(I have also used that
$dist(t,F)=|t-{1\over 3}tr(t)|_E$, in the notation of \cite{Rep},
which may not be valid if $p=3$. The modification needed for 
$p=3$ is minor and does not affect the final form of our results.)
\par
Under the isomorphism between $T_G$ and $T_H$ (see \S A.3 of the appendix),
on the $H=(SL(2)\times SL(2))/\{\pm 1\},PGL(3)$ side,
the subregular term has the form, up to a factor of $|\lambda|^{-1}$,    
$$
\cases
(|(3\alpha+2\beta)(X)|)^{-1}d_{T_H}\cdot 
\mu_{{\Cal O}_1(F)}(f^H),&\text{ for }
H=A_1\times A_1,\\
|(3\alpha+\beta)(X)(3\alpha+2\beta)(X)|^{-1}d'_{T_H}
\cdot \mu'_{{\Cal O}_1(F)}(f^H),&\text{ for }H=A_2,\\
\endcases
\tag 6.8
$$
provided $T$ is either biquadratic or cubic (respectively) and
$X$ is regular. 
Observe that, as a function of $X$, the $G_2$-side differs 
from the $H$-side by $c_{\kappa, T_G,\alpha(X)}$ times 
$D_H(\gamma_{\lambda,H})$,
in the cubic case, and by $c_{\kappa, T_G,\beta(X)}$ 
times $D_H(\gamma_{\lambda,H})$,
in the biquadratic case. 

\enddocument





\proclaim{Lemma 9.1} Let $s:=\sigma_\alpha$ and 
$r:=\sigma_\alpha\sigma_\beta$.
On $E_\alpha$, $E_{\beta}$, the generators
of the (respective) local coordinates rings on
$Z_\alpha$, $Z_\beta$ are all $\sigma$-invariant, where
$\sigma=s$ or $\sigma=r^2$.
The local coordinate $x(\delta)$, restricted to either of these
divisors, is not $\sigma$-invariant if
$$
\sigma=\sigma_\beta, \sigma_\alpha\sigma_\beta,
\sigma_\beta\sigma_\alpha, \sigma_\beta\sigma_\alpha\sigma_\beta .
$$
(This last statement does not appear to be consistent 
with \cite{H1, Lemma 1.1, \S 1} but is not needed anyway.)
\endproclaim





